%% Auto-generated LaTeX version of user_manual.md
\documentclass[11pt,a4paper]{article}
\usepackage[utf8]{inputenc}
\usepackage{fontawesome5}
\usepackage[T1]{fontenc}
\usepackage{lmodern}
\usepackage{graphicx}
\usepackage{geometry}
\usepackage{float}
\geometry{margin=1in}

% Disable paragraph indentation throughout the document
\setlength{\parindent}{0pt}

\begin{document}
\pagestyle{empty}

% Top Header: Company Info (Left) and Logo (Right)
\noindent
\begin{minipage}[c]{0.5\textwidth}
    \textbf{Mlnarik Organ, s.r.o.} \\
    Cerhenska 9 \\
    29001 Sokolec \\
    Czech Republic
\end{minipage}
\begin{minipage}[c]{0.5\textwidth}
    \raggedleft
    
\includegraphics[height=1.5cm]{logo.png} % Replace with your image file
\end{minipage}

% Place title directly below the logo (avoid starting a new page)
\vspace{3cm}

\begin{center}
    {\huge \textbf{OUTPUT--MODULE16: User Manual}}\\
    \vspace{0.3cm}
    August 2026

    \noindent Applicable to: Firmware Version 1.0

\end{center}

\vspace{1cm}

\begin{center}
    \includegraphics[width=1\textwidth]{Output-module16.png}
\end{center}

\vfill

\newpage

\section{Overview}

OUTPUT--MODULE16 receives MIDI NOTE ON/OFF messages and controls up to 16 outputs via two 16-pin IDC headers (8 outputs each). Each output is a pair of +5~V and a ground switched by an N-channel MOSFET (BSS138), suitable for driving loads such as LEDs with series resistors, relays, solenoids, or other low-power DC loads.

It supports both traditional 5-pin DIN MIDI and USB-MIDI input with automatic merging of incoming MIDI data. In addition to standard MIDI capabilities, the module introduces several advanced features such as a High-Speed Analog MIDI Bus and an easy-to-use Wi-Fi configuration interface accessible via a web browser.

\section{Features}
\begin{itemize}
    \item Board dimensions 175\texttimes{}34~mm
    \item Powered via 5~V DC (board consumption approx. 0.1~A; total depends on connected loads)
    \item 16 outputs (2 headers~\texttimes~8 outputs each) --- each output is a +5~V / FET-switched ground pair
    \item Output driver: BSS138 N-channel MOSFET (max. continuous drain current 0.22~A per output; 0.1~A recommended for reliable long-term operation)
    \item MIDI IN (5-pin DIN)
    \item MIDI OUT (5-pin DIN) with automatic merging of MIDI IN and internal messages
    \item USB MIDI input (class-compliant)
    \item Optional High-Speed Analog MIDI Bus --- supports up to 4\texttimes{} the standard MIDI baud rate for expanded bandwidth
    \item Wi-Fi Configuration Interface --- configure module settings via a web browser by connecting to the module's built-in Wi-Fi access point
\end{itemize}

\section{Introduction to MIDI Device Connectivity and Latency Considerations}

In practical applications, it is common to use multiple MIDI modules in a setup, often combining devices from various manufacturers. There are two basic approaches to connecting MIDI devices to an end device such as a computer or Digital Audio Workstation (DAW).


\begin{figure}[H]
    \centering
    \includegraphics[width=0.2\textwidth]{Analog_chaining.drawio.png}
    \caption{Chaining via the analog MIDI bus --- this typically involves using traditional 5-pin DIN connectors (though other connectors may also be used).}
\end{figure}

\begin{figure}[H]
    \centering
    \includegraphics[width=0.8\textwidth]{USB_star.drawio.png}
    \caption{Connecting in a "star" topology via USB.}
\end{figure}

\begin{center}
    \fbox{
    \begin{minipage}{0.95\linewidth}
        \faInfoCircle~\textbf{Note: Latency considerations} If you prefer chaining via the analog MIDI bus and use many MIDI devices, latency can become a significant issue. The transmission of a single standard MIDI message takes approximately 1\,ms. As MIDI data is passed through multiple devices in the chain, this latency can accumulate and become noticeable.
    \end{minipage}
    }
\end{center}

One of the key features of the OUTPUT--MODULE16 is its optional High-Speed Analog MIDI Bus mode. When activated, MIDI messages are transmitted at four times the standard MIDI speed. However, please note that this high-speed mode is not compatible with third-party MIDI devices, as they are designed to operate at the standard MIDI baud rate.

\section{Setup and installation}

\begin{figure}[H]
    \centering
    \includegraphics[width=1\textwidth]{labeled-board.drawio.png}
    \caption{Overall layout of the board}
\end{figure}

\subsection{Power}
\begin{itemize}
  \item Voltage: 5~V DC (regulated)
  \item Board self-consumption: 0.1~A max (when no outputs are active)
\end{itemize}

The total current drawn from the supply depends heavily on the loads connected to the outputs. Each output can sink up to 0.22~A (BSS138 maximum), but 0.1~A per output is recommended for reliable long-term operation. The sum of all active output currents plus the board's own 0.1~A determines the total power supply requirement.

\begin{center}
    \fbox{
    \begin{minipage}{0.95\linewidth}
        \faInfoCircle~\textbf{Power supply recommendations:}
        \begin{itemize}
            \item Use a regulated 5~V supply. Do not exceed 5.3~V. Reverse polarity can damage the device.
            \item Powering via USB is possible only when the total current (board + loads) stays within the USB source capability (typically 0.5~A for USB 2.0).
            \item For driving multiple outputs simultaneously or higher loads, use an external 5~V power supply with sufficient current capacity.
            \item When powering many modules together, use a dedicated power supply with adequate current rating or a suitable powered USB hub.
        \end{itemize}
    \end{minipage}
    }
\end{center}

\subsection{Output Wiring}

The board provides two 16-pin IDC headers (Outputs~1 and Outputs~2), each with 8 output pairs. Each pair consists of a +5~V pin and a controlled ground pin that are adjacent on the header. The ground side is switched by a BSS138 N-channel MOSFET.

To connect a load (e.g.\ an LED with a series resistor):
\begin{itemize}
    \item Take one adjacent pair from a flat ribbon cable.
    \item Connect the +5~V wire to the positive terminal of your load.
    \item Connect the controlled ground wire to the negative terminal of your load.
\end{itemize}

When a MIDI NOTE ON message is received on the configured channel and note number, the corresponding MOSFET turns on, connecting the output to ground and completing the circuit. A MIDI NOTE OFF message turns the output off.

\begin{center}
    \fbox{
    \begin{minipage}{0.95\linewidth}
        \faInfoCircle~\textbf{MIDI velocity handling:}
        \begin{itemize}
            \item The module ignores the velocity value of NOTE ON messages. Any NOTE ON with non-zero velocity turns the output on regardless of the actual velocity number.
            \item A NOTE ON message with velocity 0 is treated as NOTE OFF and turns the output off. This follows the MIDI specification's recommended practice ("Running Status" compatible).
        \end{itemize}
    \end{minipage}
    }
\end{center}

\begin{figure}[H]
    \centering
    \includegraphics[width=1\textwidth]{outputs-wiring.png}
    \caption{Output wiring diagram --- example with LEDs and series resistors}
\end{figure}

\begin{center}
    \fbox{
    \begin{minipage}{0.95\linewidth}
        \faInfoCircle~\textbf{Output specifications:}
        \begin{itemize}
            \item Each output: +5~V (fixed) + FET-switched ground (BSS138)
            \item Max. continuous current per output: 0.22~A (BSS138 rating); 0.1~A recommended for reliable long-term operation
            \item The +5~V pin is directly connected to the board's 5~V supply --- do not short it to ground.
            \item Always use a current-limiting resistor with LEDs and other low-impedance loads.
            \item The module is not designed for inductive loads (motors, large relays) without additional flyback protection.
        \end{itemize}
    \end{minipage}
    }
\end{center}

\section{Analog MIDI Connection}

The image below shows how to connect to the analog MIDI bus using DIN-5 connectors. However, it is not necessary to use connectors between the modules --- you can also use wires connected directly to the screw terminals. Shielded cables are recommended.

\begin{figure}[H]
    \centering
    \includegraphics[width=1\textwidth]{MIDI-schema.png}
    \caption{Connection analog MIDI bus}
\end{figure}


\section{Setting Up the Module - Wi-Fi Access Point Mode}

You can conveniently set up the module using any web browser on your computer or mobile phone. The module can start in a special mode where it creates a Wi-Fi access point that you can connect to.

\begin{enumerate}
    \item Unplug the module from any power source (USB and external power).
    \item Press and hold the WiFi button.
    \item Continue holding the button while powering on the module (e.g., by connecting an external power source).
    \item The Status LED will light up steadily, indicating that the module is now in Wi-Fi access
    point mode.
    \item On your computer (with a Wi-Fi adapter) or smartphone, go to the Wi-Fi settings. You
    should see a new Wi-Fi network named "Output Module16". Connect to it.

    \faInfoCircle~\textbf{Note}: While connected to this network, your device will lose access to the internet.
    \item Open a web browser (such as Google Chrome) and enter the following IP address in the address bar: 192.168.4.1 --- the fixed IP address of the MIDI module's web server. You should now see the settings page, as shown in the image below.
    \item Use the dropdown menus to configure all the necessary settings, then press the
    "Save Settings" button. If needed, you can restore the factory default values by
    pressing the "Reset to Defaults" button.
    \item Once you're done, disconnect the module from the power source. Your settings will be saved and automatically applied the next time the module starts.

    \faInfoCircle~\textbf{Note}: While in Wi-Fi access point mode, the module does not process MIDI data. MIDI communication is disabled during setup.
\end{enumerate}

\begin{figure}[H]
    \centering
    \includegraphics[width=0.8\textwidth]{app-settings-window.png}
    \caption{Settings page in Wi-Fi Access Point mode}
\end{figure}

The firmware provides a minimal HTTP-based settings page when the unit runs in Wi-Fi Access Point mode. The page allows configuring the following properties:

\begin{description}
  \item[MIDI serial bus rate] Toggle between the standard MIDI baud rate (31250) and a higher UART rate (115200). High-speed mode is not compatible with third-party MIDI equipment.
  \item[Output header 1 / Output header 2] Each header covers 8 outputs. The following settings are available per header:
    \begin{itemize}
      \item \textbf{MIDI channel} (range: 1--16). The module responds to NOTE ON/OFF messages on this channel.
      \item \textbf{Base MIDI note} (range: 0--119). The first output (pin~1) maps to this note; subsequent outputs increment the note number (output~2 = base+1, output~3 = base+2, etc.).
    \end{itemize}
\end{description}

\section{Status LED behavior}
\begin{itemize}
  \item Startup: brief blink sequence
  \item MIDI activity: short blink when MIDI bytes are received
  \item AP mode: LED remains steadily lit while the AP is active
\end{itemize}

\end{document}
